\documentclass[10pt,a4paper]{article}
\usepackage[T1]{fontenc}
\usepackage[utf8]{inputenc}
\usepackage{lmodern}
\usepackage[czech]{babel}
\usepackage[margin=1.5cm]{geometry}
\usepackage{amsmath,amssymb,mathtools}
\usepackage{xcolor,cancel,enumitem}
\usepackage{tikz,pgfplots}
\pgfplotsset{compat=1.18}
\usetikzlibrary{arrows.meta,calc,patterns}
\pagestyle{empty}
\setlength{\parindent}{0pt}
\setlength{\parskip}{4pt}
\newcommand{\vysl}[1]{\boxed{\color{red!75!black}#1}}
\newcommand{\pozn}[1]{{\color{red!75!black}#1}}
\newcommand{\dx}{\,\mathrm{d}x}
\newcommand{\dt}{\,\mathrm{d}t}
\newcommand{\subst}[1]{\left|\begin{array}{l}#1\end{array}\right|}
\DeclareMathOperator{\arctg}{arctg}
\DeclareMathOperator{\arccotg}{arccotg}
\DeclareMathOperator{\tg}{tg}
\DeclareMathOperator{\cotg}{cotg}
\usepackage{tcolorbox}
\newtcolorbox{pomocny}{colframe=red!75!black, colback=white, coltitle=white, fonttitle=\bfseries, title={Pomocný výpočet:}, sharp corners}
\newcommand{\dy}{\,\mathrm{d}y}
% znaménkové schéma: \znam{bod/popis/typ,...}{znaménka v intervalech}
% typ: o = vyloučený bod (kroužek), x = obyčejný bod
\newcommand{\znam}[2]{%
\begin{tikzpicture}[baseline=-3pt]
\def\n{0}
\foreach \p/\l/\t [count=\i] in {#1} {\xdef\n{\i}}
\pgfmathsetmacro{\L}{1.6*(\n+1)}
\draw (0,0) -- (\L,0);
\foreach \p/\l/\t [count=\i] in {#1} {
  \pgfmathsetmacro{\xx}{1.6*\i}
  \ifnum\pdfstrcmp{\t}{o}=0 \draw[fill=white] (\xx,0) circle (2.5pt); \draw (\xx-0.09,0) -- (\xx+0.09,0) (\xx,-0.09) -- (\xx,0.09);
  \else \draw (\xx,0.1) -- (\xx,-0.1); \fi
  \node[below=2pt] at (\xx,0) {\small $\l$};
}
\foreach \s [count=\i] in {#2} {
  \pgfmathsetmacro{\xx}{1.6*\i-0.8}
  \node[above] at (\xx,0) {$\s$};
}
\end{tikzpicture}}
\newcommand{\aug}[1]{\left(\begin{array}{rrr|rrr}#1\end{array}\right)}
\newcommand{\krok}[1]{\text{\small\color{orange!80!black}#1}}

\newcommand{\vv}[1]{\overrightarrow{#1}}
\newcommand{\norm}[1]{\left\|#1\right\|}

\begin{document}

\section*{Vlastní čísla, vektory -- řešení}

\textbf{\textcircled{1}} Ověřte, že $\lambda=9$ je vlastní číslo matice
$A=\begin{pmatrix} 7 & -2 & -2\\ -2 & 1 & 4\\ -2 & 4 & 1 \end{pmatrix}$
a určete odpovídající vlastní vektor.

\hrule
\medskip
{\color{blue!60!black}Stačí určit, zda $|A-\lambda E|=0$ pro $\lambda=9$:}
\[
\begin{vmatrix} 7-9 & -2 & -2\\ -2 & 1-9 & 4\\ -2 & 4 & 1-9 \end{vmatrix}
= \begin{vmatrix} -2 & -2 & -2\\ -2 & -8 & 4\\ -2 & 4 & -8 \end{vmatrix}
\ \krok{$R_2+(-1)R_1,\ R_3+(-1)R_1$}\ =
\]
\[
= \begin{vmatrix} -2 & -2 & -2\\ 0 & -6 & 6\\ 0 & 6 & -6 \end{vmatrix}
\ \krok{$R_3+R_2$}\ =
\underbrace{\begin{vmatrix} -2 & -2 & -2\\ 0 & -6 & 6\\ 0 & 0 & 0 \end{vmatrix}}_{(*)}
= \vysl{0}
\]
$\Rightarrow$ \pozn{$\lambda=9$ je vlastní číslo.}

{\color{blue!60!black}Vlastní vektor: $(A-\lambda E)\vec x=\vec o$ pro $\lambda=9$:}
\[
(*)\quad
\begin{pmatrix} -2 & -2 & -2\\ 0 & -6 & 6\\ 0 & 0 & 0 \end{pmatrix}
\begin{pmatrix} x_1\\ x_2\\ x_3 \end{pmatrix} = \begin{pmatrix} 0\\ 0\\ 0 \end{pmatrix}
\]
\begin{align*}
-2x_1-2x_2-2x_3 &= 0\\
-6x_2+6x_3 &= 0
\end{align*}
\begin{align*}
x_3 &= t\\
-6x_2+6x_3 &= 0\\
6x_2 &= 6x_3\\
x_2 = x_3 &= t\\
-2x_1-2t-2t &= 0 \ \Rightarrow\ 2x_1=-4t \ \Rightarrow\ x_1=-2t
\end{align*}
\pozn{vlastní vektor:} \quad $\vysl{\begin{pmatrix} -2t\\ t\\ t \end{pmatrix}}$

\bigskip
\hrule
\medskip

\textbf{\textcircled{2}} Jsou dány body: $A[1,2,2]$, $B[3,1,2]$, $C[5,1,1]$.

\textbf{\pozn{a) úhel $ABC$}} $=$ úhel vektorů $\vv{BA}$, $\vv{BC}$:
\begin{align*}
\vv{BA} &= A-B = (-2,1,0)\\
\vv{BC} &= C-B = (2,0,-1)\\
\cos\varphi &= \frac{\vv{BA}\cdot\vv{BC}}{\norm{\vv{BA}}\cdot\norm{\vv{BC}}}
= \frac{(-2,1,0)\cdot(2,0,-1)}{\sqrt{(-2)^2+1^2+0^2}\cdot\sqrt{2^2+0^2+(-1)^2}}
= \frac{-2\cdot2+1\cdot0+0\cdot(-1)}{\sqrt5\cdot\sqrt5} = -\frac45
\end{align*}
\[
\pozn{\sphericalangle ABC=\varphi=\arccos\left(-\frac45\right)}
\]

\textbf{\pozn{b) obsah $\triangle ABC$:}}
\[
S_\triangle = \frac12\norm{\vv{BA}\times\vv{BC}} = \frac12\sqrt{(-1)^2+(-2)^2+(-2)^2} = \frac12\cdot\sqrt9 = \vysl{\frac32}
\]
\[
\vv{BA}\times\vv{BC} = \begin{vmatrix} \vec i & \vec j & \vec k\\ -2 & 1 & 0\\ 2 & 0 & -1 \end{vmatrix}
= \vec i\cdot(-1-0)-\vec j\cdot(2-0)+\vec k\cdot(0-2) = (-1,-2,-2)
\]

\textbf{\pozn{c) délka úsečky $AB$:}}
\[
|AB| = \norm{\vv{AB}} = \norm{\vv{BA}} = \sqrt{(-2)^2+1^2+0^2} = \vysl{\sqrt5}
\]

\end{document}
